\subsection{Sparse regression and transfer to an outbreak curve}
\label{sec:track-e}

Here we compare two very different ways of discovering governing equations
from data: sparse symbolic regression and gradient-based neural fitting. We
also test separately whether the cross-domain stability fixes of
Section~\ref{sec:crossdomain} generalize beyond the mechanistic system that
produced them.

\subsubsection{Noise robustness: KAN-ODE vs.\ sparse symbolic regression}

We fit SINDy (Section~\ref{sec:sparse-method}) to the same noise-corrupted
Lotka-Volterra observations that the noise-sweep checkpoints
(Section~\ref{sec:dt-noise}) were trained on, and score both methods against
clean ground truth.

\begin{expectbox}
\expectrow{Question}{How does a sparse, symbolic method compare to a dense, 240-parameter neural method as observation noise increases?}
\expectrow{Expected}{Sparse regression's tiny library should make it more robust to noise; KAN-ODE's larger parameter count should make it more accurate.}
\expectrow{Observed}{Both expectations hold, but the KAN-ODE advantage collapses $170\times$ across the sweep, from $1110\times$ at zero noise to $6.5\times$ at $\sigma=0.1$.}
\end{expectbox}

\begin{table}[t]
\centering
\footnotesize
\renewcommand{\arraystretch}{1.1}
\caption{Extrapolation MSE vs.\ observation noise $\sigma$.}
\label{tab:sindy-noise}
\begin{tabularx}{\columnwidth}{@{}C{1.2cm} C{2.1cm} C{2.1cm} Y@{}}
\toprule
\hdrrow \thh{Noise level $\sigma$} & \thh{Sparse regression} & \thh{KAN-ODE} & \thh{KAN-ODE advantage}\\
\midrule
0.00 & $9.90\times10^{-2}$ & $\mathbf{8.92\times10^{-5}}$ & $1110\times$\\
0.01 & $1.58\times10^{-1}$ & $1.22\times10^{-3}$ & $130\times$\\
0.05 & $1.48\times10^{-1}$ & $2.54\times10^{-2}$ & $5.8\times$\\
0.10 & $8.89\times10^{-1}$ & $1.37\times10^{-1}$ & $6.5\times$\\
\bottomrule
\end{tabularx}
\end{table}

\begin{figure*}[t]
\centering
\includegraphics[width=0.85\textwidth]{figures/track_e/sindy_vs_kan_noise.png}
\caption{Training and extrapolation MSE vs.\ noise level, both methods.}
\label{fig:sindy-noise}
\end{figure*}

KAN-ODE wins on absolute accuracy at every level, but the margin collapses
$170\times$ across the sweep (sparse regression degrades $9.0\times$ from
$\sigma=0$ to $0.1$; KAN-ODE degrades $1540\times$). The six-term library of
sparse regression caps how much noise can move it. The KAN's 240 free
parameters fit the noise directly, and the error compounds through the
solver. Per-term recovery shows that the bilinear cross-term ($xy$) is
recovered to within $5.7\%$ at every noise level. This is exactly the term a
single-layer KAN edge cannot represent directly: each edge is a function of
only one input variable, and a layer sums its edges additively. Sparse
regression writes $xy$ as a literal library column, while a KAN can only
approximate it through composition across a hidden layer. A term-by-term
symbolic comparison would therefore compare a coefficient against a
distributed representation, so we compare trajectory reconstruction instead.

\begin{figure*}[t]
\centering
\includegraphics[width=0.85\textwidth]{figures/track_e/sindy_vs_kan_trajectories.png}
\caption{Reconstructed trajectories at $\sigma=0$ and $\sigma=0.1$: sparse
regression loses amplitude under noise; KAN-ODE keeps amplitude but drifts
in phase.}
\label{fig:sindy-trajectories}
\end{figure*}

\subsubsection{Edge pruning: testing whether a term-by-term comparison is even possible}

Part of the original motivation for a smooth, globally-supported basis such
as RBF is that a sparse, pruned KAN edge structure could in principle be read
back as a closed-form symbolic law. Sparse regression reaches the same
interpretability goal by construction. The base paper gets there by adding
the $L_1$ penalty of Equation~\ref{eq:l1-sparsity} to the training loss for
symbolic pruning. Every checkpoint we trained used $\lambda=0$: the plain
mean-squared-error term alone, with no sparsity pressure. A KAN-ODE analogue
of a sparse coefficient list would instead come from pruning weak edges
\emph{post hoc} and reading off what survives. We rank each edge by mean
spline-coefficient magnitude and mask the weakest, applied to the existing
(never retrained) noise-sweep checkpoints. Removing the two weakest of forty
edges raises the noise-free full-horizon MSE from $8.90\times10^{-5}$ to
$3.43$, a factor of $3.9\times10^{4}$ and worse than predicting the constant
dataset mean ($2.91$). No checkpoint was trained with sparsity-inducing
regularization, so nothing during training ever asked the network to
concentrate its fit onto a few edges. The flat plateau at $\approx3.5$--$3.8$
across all pruning levels of 25\% or more is the signature of a model where
every edge carries part of the fit. It confirms that the pruning procedure
itself works correctly on a network that has no prunable structure to find.

\begin{figure*}[t]
\centering
\includegraphics[width=0.85\textwidth]{figures/track_e/pruning_degradation.png}
\caption{Full-horizon MSE vs.\ edge-pruning fraction, all four noise levels.}
\label{fig:pruning}
\end{figure*}

\subsubsection{Testing whether SIR's fixes transfer to a non-mechanistic epidemic curve}

The outbreak curve of Section~\ref{sec:datasets}, which no compartmental ODE
generated, tests whether the SIR fixes of Section~\ref{sec:crossdomain}
generalize by mechanism rather than by coincidence of the SIR system's own
structure. The same init-time diagnostic we used to root-cause SIR
reproduces an almost identical speed-mismatch signature ($18.77\times$ vs.\
SIR's $19\times$). This justifies a time rescaling by a factor of 24
\emph{a priori}, before any training.

\begin{expectbox}
\expectrow{Question}{SIR's three stability fixes were each derived from a property of the SIR generator itself -- do they still work on data no compartmental model produced?}
\expectrow{Expected}{Only the unit-rescaling fix, which assumes nothing about the data's structure, should transfer; the two structural fixes should not.}
\expectrow{Observed}{One in three predictions was wrong: the vanishing-field gate transfers and is the best single option, despite being predicted inapplicable.}
\end{expectbox}

\begin{table*}[t]
\centering
\small
\renewcommand{\arraystretch}{1.15}
\caption{Which of SIR's three fixes transfer to a non-mechanistic curve.}
\label{tab:fix-transfer}
\begin{tabularx}{\textwidth}{@{}L{4.2cm} L{5.0cm} C{3.0cm} Y@{}}
\toprule
\hdrrow \thh{Fix} & \thh{Precondition} & \thh{Holds?} & \thh{Outcome}\\
\midrule
Time rescaling & None -- a unit change & Yes & Removes a $263{,}049\times$ gradient spike\\
Conservation via subspace projection & An exact invariant, $\sum y=c$ & No ($\sum y$ spans $0.002$--$1.54$) & Inapplicable; $25\times$ worse fit\\
Vanishing-field gate at the equilibrium point & $y_d=0$ is a manifold of equilibria & Yes (unexpectedly) & Best single option, $47\times$ better fit\\
\bottomrule
\end{tabularx}
\end{table*}

The vanishing-field gate was \emph{predicted} inapplicable: the infected
fraction at $t=0$ is $0.0027$, which seemed too small to avoid freezing the
field entirely. A control arm overturns that prediction. Gating the learned
field by the initial state correctly encodes that the rate of change is
proportional to current prevalence. It is the only configuration that
reproduces the existence of a peak (day 48 predicted vs.\ true day 45) from a
training window that never observed one.

\begin{figure*}[t]
\centering
\includegraphics[width=0.85\textwidth]{figures/track_e/epidemic_time_scale_sweep.png}
\caption{Time-rescaling sweep: accuracy and gradient-spike ratio vs.\ the
rescaling factor.}
\label{fig:time-scale-sweep}
\end{figure*}

\paragraph{The extrapolation runaway traces to the train/test split.} The
default split lands exactly on the infected peak (0\% of the decay phase is
in the training window), so every model is asked to predict a monotone decay
it has never observed. Moving the split just 15 days later collapses the
predicted maximum from $6.18\times$ the true peak to $0.964$ (true maximum
$=1.0$), turning over at day 47 against a true day 45. The runaway is a
property of where the window ends, not of the method fitting the data.

\begin{figure*}[t]
\centering
\includegraphics[width=0.85\textwidth]{figures/track_e/real_epidemic_fit.png}
\caption{Epidemic-curve fit: default split (runaway), the vanishing-field-gate
arm, and the fifteen-day-later split control.}
\label{fig:real-epidemic-fit}
\end{figure*}

\begin{keypoint}[Verdict]
Time rescaling transfers unconditionally because it is a unit change with no
structural precondition. The other two fixes are structural priors that
transfer exactly insofar as their structure is actually present in the new
data: the conservation projection's precondition is confirmable in advance
by directly measuring the compartment sum, while the vanishing-field gate's
precondition turned out to require running the control arm to settle --
argument alone predicted the wrong outcome. That arm cost 27 minutes and
overturned the prediction.
\end{keypoint}
